% To compile using latexmk on the command line, run the following: 
% latexmk -pdf specificAims.tex

\documentclass[11pt]{article}

% ---------- Page formatting ----------
\usepackage{geometry}
\geometry{margin=1in}

\usepackage{lmodern}
\usepackage[T1]{fontenc}
\usepackage{setspace}
\onehalfspacing

% ---------- Graphics ----------
\usepackage{graphicx}
\usepackage{float}

% ---------- Math ----------
\usepackage{amsmath}
\usepackage{amssymb}
\usepackage{bm}

% ---------- ACS citations ----------
% \usepackage{achemso}
% \setkeys{acs}{articletitle=true}

% ---------- Title ----------
\title{\textbf{Computational Identification of Unknown Organic Contaminants and Molecules in Remote Environments via High-Resolution Mass Spectrometry}}
\author{Raelyn Brooks}
\date{\today}

\begin{document}

\maketitle

\nocite{*}

\begin{doublespace}

\section{Goal}
    This project aims to characterize the chemical footprint of global pollution in remote high-altitude environments by identifying unknown organic contaminants in snowmelt and river water using high-resolution mass spectrometry (HRMS).
    I hypothesize that long-range atmospheric transport leads to the accumulation of diverse and previously uncharacterized organic molecules in these regions.
    By combining non-targeted analysis with molecular networking, this study will reveal both the identity and distribution patterns of these contaminants.

\section{Specific Aims}
    \subsection{Aim 1: Identify known and suspected contaminants using database-driven screening}
        \begin{itemize}
            \item {\textbf{What (Aim):}} Identify known and suspected organic contaminants in high-altitude water samples through HRMS-based suspect screening against curated chemical databases.
            \item {\textbf{How (Approach):}} Mass spectral features will be matched against databases (e.g., environmental contaminant libraries) using accurate mass and fragmentation patterns to annotate compounds with varying confidence levels.
            \item {\textbf{Impact:}} This aim will establish a baseline of identifiable contaminants present in remote environments, linking detected compounds to known anthropogenic sources.
        \end{itemize}

    \subsection{Aim 2: Discover structurally related and unknown compounds using molecular networking}
        \begin{itemize}
            \item {\textbf{What (Aim):}} Characterize unknown and structurally related molecules by organizing HRMS data into molecular networks.
            \item {\textbf{How (Approach):}} MS/MS spectral similarity will be used to construct molecular networks, grouping compounds into clusters that reveal structural relationships and enable annotation of unknowns based on known neighbors.
            \item {\textbf{Impact:}} This aim will expand detection beyond known compounds, enabling discovery of novel or transformation products and improving chemical annotation coverage.
        \end{itemize}

    \subsection{Aim 3: Identify patterns in contaminant distribution using unsupervised machine learning}
        \begin{itemize}
            \item {\textbf{What (Aim):}} Determine patterns and relationships in contaminant profiles across samples using unsupervised machine learning techniques.
            \item {\textbf{How (Approach):}} Techniques such as hierarchical clustering and principal component analysis (PCA) will be applied to HRMS data to group samples and identify trends based on chemical composition.
            \item {\textbf{Impact:}} This aim will reveal underlying structure in the data, helping to infer sources, transport pathways, and environmental factors influencing contaminant distribution.
        \end{itemize}

%\begin{figure}[H] % Suggests figure placement (here, top, bottom)
%    \centering
%    \includegraphics[scale=0.50]{SMILEStoMolecule.png}
%    \caption{Visuals of SMILES notation to actual molecules}
%    \label{fig:SMILEStomolecule} % Unique label for referencing
%\end{figure}

\end{doublespace}

% ---------- References ----------
\newpage
\bibliographystyle{plain}
\bibliography{bibliography.bib}

\end{document}
